\section*{Заключение}
\addcontentsline{toc}{section}{Заключение}

В ходе выполнения курсового проекта было разработано веб-приложение для автоматизации учета книг и читателей библиотеки. Цель работы заключалась в создании системы, позволяющей вести каталог книг, учитывать читателей, оформлять выдачу и возврат книг, просматривать статистику и выполнять выгрузку данных. Поставленная цель была достигнута.

На первом этапе была рассмотрена предметная область и определены основные задачи, которые должно решать приложение. Были выделены основные группы пользователей: читатель, библиотекарь и администратор. Для каждой группы определены доступные функции и ограничения. Читатель может просматривать каталог и карточки книг, библиотекарь выполняет основные операции учета, а администратор управляет учетными записями и ролями пользователей.

На этапе проектирования были сформулированы функциональные требования к системе, построены диаграммы вариантов использования, активности, классов и последовательности, а также разработана ER-диаграмма базы данных. В модели данных были выделены сущности книги, автора, категории, читателя и выдачи. Для поддержки нескольких авторов и категорий у одной книги были использованы промежуточные сущности связи. Также была предусмотрена загрузка обложек книг и хранение сведений о составе томов конкретного издания.

В процессе разработки было реализовано веб-приложение на базе фреймворка Django. В приложении настроены маршруты, модели, формы, HTML-шаблоны, пользовательские стили и работа с медиафайлами. Для хранения данных используется SQLite, что является достаточным решением для учебного проекта и локальной демонстрации работы приложения.

В приложении реализованы основные функции учета библиотечного фонда: просмотр каталога книг, поиск, фильтрация, сортировка, раскрытие разных изданий одной книги, просмотр карточки книги, добавление, редактирование и удаление записей. Для карточки книги предусмотрены отображение обложки, авторов, категорий, года издания, состава томов, описания и доступности экземпляров.

Также реализован учет читателей. Библиотекарь может просматривать список читателей, искать читателя по ФИО, телефону, электронной почте или логину, открывать карточку читателя и редактировать его данные. Для номеров телефонов реализована нормализация, благодаря чему данные сохраняются в едином формате.

Для учета движения книг реализованы оформление выдачи и возврат книги. При создании выдачи система проверяет наличие доступных экземпляров и корректность дат. После оформления выдачи количество доступных экземпляров уменьшается, а после возврата увеличивается. Журнал выдач позволяет просматривать записи, выполнять поиск, фильтрацию по статусу, сортировку и выгрузку данных в CSV.

Отдельно реализован раздел статистики. В нем отображаются показатели по выбранному периоду: количество читателей с книгами, активные выдачи, общее число выдач, просроченные выдачи, состояние фонда, популярные категории, последние выдачи и график выдач по месяцам. По умолчанию статистика формируется за период с первого дня текущего месяца по текущую дату.

В приложении настроено разграничение доступа по ролям. Читатель не имеет доступа к журналу выдач, статистике и функциям управления книгами. Библиотекарь работает с основным интерфейсом приложения. Администратор управляет пользователями и группами через административную панель Django. Это позволяет разделить обязанности пользователей и ограничить доступ к служебным функциям.

В результате выполнения работы было создано функциональное веб-приложение, соответствующее поставленным требованиям. Разработанная система может использоваться как учебный пример автоматизации работы небольшой библиотеки. При дальнейшем развитии проекта можно добавить расширенный импорт данных из CSV, уведомления читателей о приближении срока возврата, личный кабинет читателя, историю изменений записей и более детальную аналитику по фонду.